% Path: docs/report_skeleton.tex
\documentclass[11pt,a4paper]{article}

\usepackage[utf8]{inputenc}
\usepackage[T1]{fontenc}
\usepackage{geometry}
\usepackage{booktabs}
\usepackage{graphicx}
\usepackage{float}
\usepackage{hyperref}
\usepackage{longtable}
\usepackage{array}
\usepackage{enumitem}

\geometry{margin=1in}

\title{CNN-LSTM Image Captioning on Flickr8k}
\author{Giuseppe Dimonte}
\date{\today}

\begin{document}

\maketitle

\begin{abstract}
This report presents the implementation and evaluation of an image captioning system based on a convolutional encoder and a recurrent decoder. The project follows the requirements of the Deep Learning and Generative Models assignment on the Flickr8k dataset. The final pipeline includes caption preprocessing, vocabulary construction, supervised CNN-LSTM training, BLEU-based evaluation, and qualitative visualization. A controlled experimental comparison was conducted on four configurations: a baseline setup, an ablation without scheduled sampling, an ablation without label smoothing, and a fine-tuning configuration. The best result was achieved by the fine-tuned model, which obtained a BLEU-4 score of 0.0649, outperforming the baseline BLEU-4 score of 0.0472.
\end{abstract}

\section{Introduction}

Image captioning is a multimodal task that combines computer vision and natural language generation. Given an input image, the objective is to produce a textual description that captures the relevant visual content. This task is naturally formulated as an image-to-sequence problem: an image is first encoded into a visual representation, then a language model generates a sequence of words conditioned on that representation.

The assignment required the implementation of a CNN-LSTM architecture on Flickr8k. The chosen approach uses a pretrained convolutional neural network to extract visual cues and an LSTM decoder to generate the caption token by token. This design follows the classical encoder-decoder formulation for caption generation and remains suitable for an academic study where the full training and evaluation pipeline must be understood and documented end-to-end.

The main goals of the project were:
\begin{itemize}[leftmargin=*]
    \item to build a complete image captioning pipeline from data preprocessing to evaluation;
    \item to make the workflow reproducible through explicit preparation scripts and configuration files;
    \item to compare a baseline setup against targeted ablations;
    \item to collect quantitative and qualitative evidence for the final academic discussion.
\end{itemize}

\section{Assignment Requirements}

The original assignment specified the following key points:
\begin{itemize}[leftmargin=*]
    \item use the Flickr8k dataset;
    \item implement a CNN-LSTM architecture;
    \item clean the caption text by converting to lowercase and removing punctuation, numbers, and single-letter words;
    \item create a vocabulary file suitable for sequence modeling;
    \item pair image features with tokenized captions during training;
    \item use the decoder output to predict caption words at test time.
\end{itemize}

The final repository satisfies these requirements through a modular structure with separate files for preprocessing, dataset management, model definition, training, evaluation, and visualization.

\section{Dataset and Preprocessing}

\subsection{Dataset}

The project uses Flickr8k, a dataset of images paired with natural language descriptions. Each image is associated with five reference captions. For the final prepared version of the project, the following derived dataset sizes were obtained:
\begin{itemize}[leftmargin=*]
    \item 30000 training caption rows;
    \item 5000 validation caption rows;
    \item 1000 test images;
    \item 5 reference captions per test image.
\end{itemize}

\subsection{Text Cleaning}

The caption preprocessing stage applies the following transformations:
\begin{itemize}[leftmargin=*]
    \item lowercase conversion;
    \item punctuation removal;
    \item digit removal;
    \item removal of single-letter words;
    \item whitespace normalization.
\end{itemize}

This preprocessing was implemented in the project code and executed before training. The vocabulary was built from the cleaned Flickr8k captions with a minimum frequency threshold, and the final vocabulary size was 2985 tokens, including the special symbols \texttt{<pad>}, \texttt{<bos>}, \texttt{<eos>}, and \texttt{<unk>}.

\subsection{Generated Artifacts}

The data preparation step produced the following reusable files:
\begin{itemize}[leftmargin=*]
    \item \texttt{data/vocab.json}
    \item \texttt{data/Flickr8k\_text/train.csv}
    \item \texttt{data/Flickr8k\_text/val.csv}
    \item \texttt{data/Flickr8k\_text/test\_captions.json}
\end{itemize}

These files make the full training and evaluation pipeline reproducible.

\section{Model Architecture}

The implemented model follows a standard encoder-decoder design.

\subsection{Encoder}

The encoder is based on pretrained InceptionV3. The original classification head is removed, and the extracted visual features are projected into the embedding space used by the decoder. This allows the model to transform the image into a compact visual representation suitable for sequence generation.

\subsection{Decoder}

The decoder is an LSTM network that generates a caption one token at a time. During training, the decoder receives the image-conditioned representation and the tokenized caption sequence. A final linear layer maps each hidden state to vocabulary logits, allowing the model to predict the next word in the caption.

\subsection{Training Features}

The training code supports:
\begin{itemize}[leftmargin=*]
    \item scheduled sampling;
    \item label smoothing;
    \item gradient clipping;
    \item checkpoint saving;
    \item checkpoint resume;
    \item BLEU-based evaluation with 5 reference captions per image.
\end{itemize}

\section{Implementation Details}

The project was organized into modular components:
\begin{itemize}[leftmargin=*]
    \item \texttt{src/vocab\_builder.py}: caption cleaning, vocabulary building, split export;
    \item \texttt{src/prepare\_data.py}: convenience wrapper for data preparation;
    \item \texttt{src/dataset.py}: dataset loading and batching;
    \item \texttt{src/model.py}: encoder-decoder model;
    \item \texttt{src/train.py}: training loop and checkpoint logic;
    \item \texttt{src/eval.py}: caption generation and BLEU computation;
    \item \texttt{src/visualize.py}: qualitative figure generation;
    \item \texttt{src/ui\_streamlit.py}: optional interactive demo.
\end{itemize}

During the revision phase, several issues were fixed to make the project executable and defensible:
\begin{itemize}[leftmargin=*]
    \item the Flickr8k caption parser was corrected to match the real tab-separated format;
    \item missing train/validation/test artifacts were generated automatically;
    \item loss alignment was corrected for next-token prediction;
    \item encoder fine-tuning was made effective when enabled;
    \item the training and evaluation paths were aligned with the actual repository structure.
\end{itemize}

\section{Experimental Setup}

\subsection{Environment}

All experiments were executed in the \texttt{captioning} Python environment with:
\begin{itemize}[leftmargin=*]
    \item Python executable: \texttt{C:\textbackslash Users\textbackslash Giuseppe\textbackslash anaconda3\textbackslash envs\textbackslash captioning\textbackslash python.exe}
    \item PyTorch version: \texttt{2.9.1+cpu}
    \item compute device: CPU
\end{itemize}

This constraint is important because it directly affects training time and explains the need for controlled experiment design.

\subsection{Compared Configurations}

Four configurations were evaluated:
\begin{enumerate}[leftmargin=*]
    \item \textbf{Baseline}: scheduled sampling and label smoothing enabled;
    \item \textbf{No Scheduled Sampling}: scheduled sampling disabled;
    \item \textbf{No Label Smoothing}: label smoothing disabled;
    \item \textbf{Fine-tuning}: CNN fine-tuning enabled.
\end{enumerate}

In the final comparison, all four configurations were run for 5 epochs.

\section{Quantitative Results}

\subsection{BLEU Comparison}

\begin{table}[H]
\centering
\caption{BLEU results for the final experimental comparison.}
\begin{tabular}{lcccc}
\toprule
Experiment & BLEU-1 & BLEU-2 & BLEU-3 & BLEU-4 \\
\midrule
Baseline & 0.2553 & 0.1556 & 0.0928 & 0.0472 \\
No Scheduled Sampling & 0.2232 & 0.1242 & 0.0717 & 0.0396 \\
No Label Smoothing & 0.2404 & 0.1435 & 0.0809 & 0.0394 \\
Fine-tuning & 0.3301 & 0.2022 & 0.1189 & 0.0649 \\
\bottomrule
\end{tabular}
\end{table}

\subsection{Training Summary}

\begin{table}[H]
\centering
\caption{Training summary for the four experiments.}
\begin{tabular}{lccc}
\toprule
Experiment & Best Val Loss & Final Train Loss & Final Val Loss \\
\midrule
Baseline & 4.1927 & 4.8016 & 4.2332 \\
No Scheduled Sampling & 4.0106 & 3.5373 & 4.0234 \\
No Label Smoothing & 3.5029 & 4.0860 & 3.5738 \\
Fine-tuning & 3.6614 & 4.7512 & 4.2861 \\
\bottomrule
\end{tabular}
\end{table}

\section{Discussion}

\subsection{Effect of Scheduled Sampling}

The removal of scheduled sampling lowers all BLEU metrics relative to the baseline. Interestingly, the training loss becomes lower without scheduled sampling, but this does not translate into better caption quality. This suggests that scheduled sampling improves robustness and generalization during generation.

\subsection{Effect of Label Smoothing}

Removing label smoothing also reduces BLEU performance relative to the baseline. This indicates that label smoothing provides a regularization benefit in the current setup, improving the quality of generated captions even when the validation loss alone might seem competitive.

\subsection{Effect of Fine-tuning}

The fine-tuning configuration achieves the best overall BLEU scores:
\begin{itemize}[leftmargin=*]
    \item BLEU-1: 0.3301
    \item BLEU-2: 0.2022
    \item BLEU-3: 0.1189
    \item BLEU-4: 0.0649
\end{itemize}

This result shows that allowing the encoder to adapt to the captioning task yields the strongest quantitative performance in this study.

\section{Qualitative Results}

Qualitative outputs were generated for all four configurations and saved in:
\begin{itemize}[leftmargin=*]
    \item \texttt{figures/baseline\_cpu/}
    \item \texttt{figures/ablation\_no\_ss\_cpu/}
    \item \texttt{figures/ablation\_no\_ls\_cpu/}
    \item \texttt{figures/ablation\_finetune\_cpu/}
\end{itemize}

For the final report version, it is recommended to include:
\begin{itemize}[leftmargin=*]
    \item one representative baseline example;
    \item one strong fine-tuned example;
    \item one failure case for discussion.
\end{itemize}

\section{Limitations}

\begin{itemize}[leftmargin=*]
    \item All experiments were executed on CPU, which significantly increased training time.
    \item Flickr8k is relatively small for image captioning, limiting the absolute quality achievable by the model.
    \item The work focuses on a classical CNN-LSTM setup and does not compare against more modern transformer-based architectures.
\end{itemize}

\section{Conclusion}

This project implemented a complete CNN-LSTM image captioning pipeline on Flickr8k, including preprocessing, training, evaluation, and qualitative visualization. The final experimental comparison shows that both scheduled sampling and label smoothing improve performance relative to their corresponding ablations. The best configuration is the fine-tuned model, which achieved the highest BLEU scores across all tested setups. The repository is now in a condition suitable for final academic reporting.

\section*{Appendix: Available Artifacts}

The following files can be cited or reused in the final documentation:
\begin{itemize}[leftmargin=*]
    \item \texttt{results/metrics\_summary.csv}
    \item \texttt{results/experiment\_summary.csv}
    \item \texttt{results/*\_predictions.json}
    \item \texttt{logs/*}
    \item \texttt{figures/*}
    \item \texttt{LAB\_LOG.md}
\end{itemize}

\end{document}
